\documentclass[11pt]{article}
\usepackage[margin=1in]{geometry}
\usepackage{amsmath,amssymb,mathtools,bm}

\title{Janiform Model Math Summary}
\author{}
\date{}

\begin{document}
    \maketitle

    \section*{11. Anchor–Deviation Extension (Mean-Reverting Layer)}

    \subsection*{11.1. Motivation}
    To allow the model to express assets with long-horizon anchoring behavior
    (e.g., FX, commodities, spreads), we introduce a minimal mean-reverting
    extension that preserves the existing Janiform core as the innovation driver.

    The goal is to add an equilibrium-restoring mechanism without modifying
    the Hawkes--NIG--news structure or introducing strong identifiability conflicts.

    \subsection*{11.2. Anchored Log-Price Dynamics}

    Let \(X_t^{\mathrm{JF}}\) denote the original Janiform return process:
    \[
        dX_t^{\mathrm{JF}} = d\big(L^+_{A_t^+} - L^-_{A_t^-} + M_t + G_t\big).
    \]

    Define the observed log-price \(X_t\) via:
    \[
        dX_t = -\kappa\big(X_t - A_t\big)\,dt + dX_t^{\mathrm{JF}},
    \]
    where:
    \begin{itemize}
        \item \(\kappa \ge 0\) is the mean-reversion speed,
        \item \(A_t\) is an anchor (equilibrium) process.
    \end{itemize}

    This nests the original model:
    \[
        \kappa = 0 \quad \Rightarrow \quad dX_t = dX_t^{\mathrm{JF}}.
    \]

    \subsection*{11.3. Anchor Process \(A_t\)}

    The anchor \(A_t\) is intentionally kept low-dimensional to preserve identifiability.
    Typical specifications include:

    \paragraph{(i) Constant anchor}
    \[
        A_t \equiv m.
    \]

    \paragraph{(ii) Deterministic seasonal anchor}
    \[
        A_t = f_{\mathrm{seasonal}}(t),
    \]
    where \(f_{\mathrm{seasonal}}\) may depend on time-of-day, day-of-week,
    day-of-year, or known calendar structure.

    \paragraph{(iii) Slowly varying anchor}
    \[
        dA_t = \mu_A(t)\,dt,
    \]
    with \(\mu_A(t)\) constrained to be low-frequency.

    In all cases, \(A_t\) is either fixed or externally estimated in the first stage,
    and not fully free during initial calibration.

    \subsection*{11.4. Deviation Formulation}

    Define the deviation from anchor:
    \[
        D_t := X_t - A_t.
    \]

    Then:
    \[
        dD_t = -\kappa D_t\,dt + dX_t^{\mathrm{JF}} - dA_t.
    \]

    This shows explicitly that:
    \begin{itemize}
        \item the Janiform process acts as the innovation driver,
        \item the deviation \(D_t\) is mean reverting toward zero,
        \item the anchor \(A_t\) shifts the equilibrium level.
    \end{itemize}

    \subsection*{11.5. Integrated Form}

    For a time interval \([t, t+\Delta]\), the solution is:
    \[
        X_{t+\Delta}
        =
        A_{t+\Delta}
        +
        e^{-\kappa \Delta}\big(X_t - A_t\big)
        +
        \int_t^{t+\Delta} e^{-\kappa(t+\Delta - s)}\,dX_s^{\mathrm{JF}}.
    \]

    This representation highlights:
    \begin{itemize}
        \item exponential decay of deviations from the anchor,
        \item filtering of Janiform increments through the kernel \(e^{-\kappa(\cdot)}\).
    \end{itemize}

    \subsection*{11.6. Interpretation}

    \begin{itemize}
        \item The Janiform core \(X_t^{\mathrm{JF}}\) continues to model:
        \begin{itemize}
            \item heavy tails,
            \item jump behavior,
            \item self-excitation,
            \item news-driven shocks,
            \item short- to medium-horizon dynamics.
        \end{itemize}

        \item The new term \(-\kappa(X_t - A_t)\) introduces:
        \begin{itemize}
            \item long-horizon anchoring,
            \item equilibrium-restoring dynamics,
            \item control over unrealistic drift at large time scales.
        \end{itemize}
    \end{itemize}

    \subsection*{11.7. Calibration Strategy (Staged)}

    To preserve stability and identifiability:

    \begin{enumerate}
        \item Calibrate the original Janiform model (\(\kappa=0\)) as before.
        \item Specify or estimate \(A_t\) (deterministic or low-frequency).
        \item Estimate \(\kappa\) using lower-frequency or aggregated data.
        \item Optionally perform limited joint refinement (though this step should be left unimplemented unless shown to be necessary).
    \end{enumerate}

    \subsection*{11.8. Identifiability Constraints}

    To minimize parameter competition:

    \begin{itemize}
        \item \(A_t\) must be low-dimensional or externally constrained.
        \item \(\kappa\) should correspond to slower time scales than Hawkes decay.
        \item No feedback from \(X_t\) or \(D_t\) into:
        \begin{itemize}
            \item Hawkes intensities,
            \item Lévy drivers,
            \item news processes.
        \end{itemize}
        \item No additional jump structure is introduced in the deviation layer.
    \end{itemize}

    \subsection*{11.9. Summary}

    This extension introduces mean reversion as a separate structural layer,
    without modifying the internal Janiform dynamics. It provides:

    \begin{itemize}
        \item a nested generalization of the original model,
        \item minimal disruption to existing calibration,
        \item explicit control over long-horizon behavior,
        \item flexibility across asset classes via the choice of \(A_t\).
    \end{itemize}

\end{document}
